\scanpage{140}
\mspage{V.12}
% archive sheet number: 139
% manuscript heading
3. \uline{Application du th. de Lefschetz \deleted{\illegible{1}} \inserted{(de la section hyperplane)} \unsure{global}\deleted{\unsure{e}}.}

% manuscript lines 1-5
\deleted{\unsure{3.1|3.4}} 3.1. Supposons que la \uline{profondeur étale rectifiée} \inserted{pour $\ell$} de $X_s$ et $X_{\unsure{\eta|y}}$ soit $\geq m$, i.e. que pour tout $x\in X_s$, on ait
\[
\overbrace{\operatorname{prof\,et}_x X_s+\dim\overline{\{x\}}}^{\inserted{\operatorname{prof\,et}_x^{*}(X_s)}}\geq m
\tag*{(3.1.1)}
\]
et de même pour $\forall\ x\in X_\eta$,
\[
\underbrace{\operatorname{prof\,et}_x(X_\eta)+\dim\overline{\{x\}}}_{\inserted{\operatorname{prof\,et}_x^{*}(X_\eta)}}\geq m,
\tag*{(3.1.2)}
\]
% manuscript lines 6-9
où ici $\overline{\{x\}}$ désigne l'adhérence dans $X_\eta$. (On note que l'hypothèse faite sur $X_\eta$ ne dépend en fait que de la fibre \uline{géométrique} $X_{\bar\eta}$).
% slanted
\margin{! Cf SGA 1 \uline{XIV} --- pour la notion de \uline{profondeur étale}.}
\deleted{Alors, \unsure{en}} Soit $d$ comme dans 2.5, et supposons \inserted{d'abord} $X$ projectif et \uline{plat sur} $S$.
% manuscript lines 10-15
Choisissons $Y$ comme dans 2.7, alors le th. de Lefschetz \deleted{\unsure{p}} \inserted{global} sous la forme SGA 1 \uline{XIV} 4.6 \uline{entraîne} \uline{[sous réserve de disposer de la résolution des singularités sur \deleted{\unsure{tous}} les \unsure{clôtures|clotures} alg. de $k$ et $K$.]} :
\[
H^i(X_s,\Lambda)\mathbin{-}H^i(Y_s,\Lambda)
\quad(\text{resp. }H^i(X_{\bar\eta},\Lambda)\mathbin{-}H^i(Y_{\bar\eta},\Lambda))
\]
est bijectif pour $i<m-d-1$, injectif pour $i=m-d-1$. Le diagramme commutatif
\diagram{\begin{tikzcd}[ampersand replacement=\&]
H^i(X_s,\Lambda)\arrow[r,"\varphi_X^i"]\arrow[d]\&H^i(X_{\bar\eta},\Lambda)\arrow[d]\\
H^i(Y_s,\Lambda)\arrow[r,"\varphi_Y^i","\sim"']\&H^i(Y_{\bar\eta},\Lambda)
\end{tikzcd}}
% manuscript lines 16-19
(où $\varphi_Y$ est bijectif en vertu de SGA 4 XVI ---) montre que \deleted{$\varphi_X^i$ est bijectif}
\[
\deleted{\unsure{\text{(3.1)}}}\qquad\varphi_X^i:H^i(X_s,\Lambda)\mathbin{-}H^i(X_{\bar\eta},\Lambda)
\]
\uline{est bijectif pour $i<m-d-1$, injectif pour} $i=m-d-1$,
